\documentclass[11pt,a4paper]{article}
\usepackage[utf8]{inputenc}
\usepackage[T1]{fontenc}
\usepackage[spanish,es-minimal,es-noshorthands]{babel}
\AtBeginDocument{\renewcommand{\sin}{\operatorname{sen}}\renewcommand{\arcsin}{\operatorname{arcsen}}}
\usepackage{lmodern}
\usepackage[sfdefault,lining]{FiraSans}
\usepackage[a4paper,top=2.3cm,bottom=2.2cm,left=1.8cm,right=1.8cm,headheight=1.3cm,headsep=0.4cm,footskip=1cm]{geometry}
\usepackage{amsmath,amssymb}
\usepackage{icomma}
\usepackage{graphicx}
\usepackage[table]{xcolor}
\usepackage{tikz}
\usepackage{fontawesome5}
\usepackage{ccicons}
\usepackage{fancyhdr}
\usepackage{lastpage}
\usepackage{tasks}
\usepackage{needspace}
\usepackage[most]{tcolorbox}
\usepackage{array}
\usepackage[hidelinks]{hyperref}
\AtBeginDocument{\hypersetup{pdftitle={Ficha de ejercicios de la unidad 3: Polinomios. 4º ESO},pdfsubject={Matemáticas A 4º ESO},pdfkeywords={polinomios, monomios, valor numérico, suma de polinomios, resta de polinomios, producto de polinomios, identidades notables, división de polinomios, regla de Ruffini, teorema del resto, raíces, factorización, factor común}}}

\definecolor{navy}{HTML}{1B2748}
\definecolor{accent}{HTML}{E8603A}
\definecolor{teal}{HTML}{1F9E90}
\definecolor{soft}{HTML}{EEF1F8}
\definecolor{muted}{HTML}{5D6685}

\setlength{\parindent}{0pt}
\setlength{\parskip}{3pt}
\settasks{label=\alph*), label-format=\bfseries\color{navy}, label-width=1.3em, item-indent=1.8em, column-sep=0.6em, after-item-skip=5pt, before-skip=3pt, after-skip=4pt}

% ---------- capçalera i peu ----------
\newcommand{\llicencia}{\href{https://creativecommons.org/licenses/by-nc-sa/4.0/deed.es}{\ccbyncsa}\ \ Licencia Creative Commons Atribución-NoComercial-CompartirIgual 4.0 Internacional (CC BY-NC-SA 4.0)}
\pagestyle{fancy}
\fancyhf{}
\fancyhead[L]{\raisebox{-0.35\height}{\includegraphics[height=0.95cm]{logo-gva.png}}}
\fancyhead[C]{\small\color{muted}Ficha de ejercicios · Unidad 3: Polinomios · 4\textsuperscript{o} ESO A}
\fancyhead[R]{\raisebox{-0.35\height}{\includegraphics[height=0.95cm]{logo-vermellar.png}}}
\fancyfoot[L]{\scriptsize\color{muted}\llicencia}
\fancyfoot[R]{\small\color{muted}\thepage}
\renewcommand{\headrulewidth}{0.4pt}
\renewcommand{\headrule}{\color{muted!40}\hrule height \headrulewidth}
\fancypagestyle{primera}{\fancyhf{}\renewcommand{\headrulewidth}{0pt}
  \fancyfoot[L]{\scriptsize\color{muted}\llicencia}
  \fancyfoot[R]{\small\color{muted}\thepage}}

% ---------- nivells i exercicis ----------
\newcommand{\estrelles}[1]{{\color{accent}\ifcase#1\or\faStar\faStar[regular]\faStar[regular]\or\faStar\faStar\faStar[regular]\or\faStar\faStar\faStar\fi}}
\newcounter{exer}
\newcommand{\ex}[1]{\par\needspace{4\baselineskip}\medskip\refstepcounter{exer}\noindent{\color{navy}\bfseries\theexer.}\ \ {\small\estrelles{#1}}\ \ }

\newcommand{\seccio}[2]{\par\needspace{16\baselineskip}\bigskip
  \begin{tcolorbox}[colback=navy,colframe=navy,coltext=white,arc=2mm,boxrule=0pt,left=3mm,right=3mm,top=1.2mm,bottom=1.2mm]
  {\large\bfseries #1\quad #2}
  \end{tcolorbox}\vspace{-2pt}}

\newcommand{\exf}[1]{\par\medskip\refstepcounter{exer}\noindent{\color{navy}\bfseries\theexer.}\ \ {\small\estrelles{#1}}\ \ }
\newcolumntype{C}[1]{>{\centering\arraybackslash}p{#1}}
% quadre de recordatori (només a la fitxa d'exercicis)
\newcommand{\recorda}[1]{\par\smallskip\begin{tcolorbox}[colback=accent!7,colframe=accent!55,arc=2mm,boxrule=0.5pt,left=2.5mm,right=2.5mm,top=1mm,bottom=1mm]\small{\color{accent}\faLightbulb[regular]}\ \textbf{Recuerda:}\ #1\end{tcolorbox}}
\newcommand{\res}[1]{#1}\newcommand{\rest}[1]{#1}

\def\fmneed{11}% que el títol d'una secció no quede sol al final del full
% =====================================================================
% Capçalera FentMates de les fitxes: un color per a cada curs.
% Els colors són els de la pàgina inicial (LEVELS a assets/js/materials.js);
% el logo ve de _eines/logo/pdf/ (python3 _eines/logo/genera.py).
%
% Ús, al final del preàmbul de la unitat:
%   % =====================================================================
% Capçalera FentMates de les fitxes: un color per a cada curs.
% Els colors són els de la pàgina inicial (LEVELS a assets/js/materials.js);
% el logo ve de _eines/logo/pdf/ (python3 _eines/logo/genera.py).
%
% Ús, al final del preàmbul de la unitat:
%   % =====================================================================
% Capçalera FentMates de les fitxes: un color per a cada curs.
% Els colors són els de la pàgina inicial (LEVELS a assets/js/materials.js);
% el logo ve de _eines/logo/pdf/ (python3 _eines/logo/genera.py).
%
% Ús, al final del preàmbul de la unitat:
%   \input{../comu/fentmates.tex}
%   \fmcurs{1eso}
% i, al començament del document:
%   \capcalera{Matemàtiques}{1r ESO}{Fitxa d'exercicis · Unitat 2}{Divisibilitat}
% Opcional abans de \input: \def\fmneed{14} (espai mínim abans de cada secció).
% =====================================================================
\makeatletter
\def\fm@logodir{../../logo/pdf}
\providecommand{\fmneed}{10}
\def\fm@cursnom{}

% clau del curs, color i color del text damunt del color
\newcommand{\fm@curs}[3]{\@namedef{fm@color@#1}{#2}\@namedef{fm@text@#1}{#3}}
\fm@curs{1eso}{E6B84A}{navy}
\fm@curs{2eso}{3B82F6}{white}
\fm@curs{3eso}{FF7A50}{white}
\fm@curs{4eso-a}{A78BFA}{white}
\fm@curs{4eso-b}{6D28D9}{white}
\fm@curs{1batx-mi}{9B1C1C}{white}
\fm@curs{1batx-ccss}{EF6B6B}{white}
\fm@curs{2batx-mii}{1E6B47}{white}
\fm@curs{2batx-ccss}{2BB3A3}{white}

\newcommand{\fmcurs}[1]{%
  \@ifundefined{fm@color@#1}{\PackageError{fentmates}{Curs desconegut: #1}{}}{}%
  \edef\fm@tmp{\noexpand\definecolor{curs}{HTML}{\@nameuse{fm@color@#1}}%
               \noexpand\colorlet{curstext}{\@nameuse{fm@text@#1}}}\fm@tmp
  \colorlet{cursfosc}{curs!72!black}%
  \def\fm@logo{\fm@logodir/vertical-#1.pdf}%
  \def\fm@marca{\fm@logodir/marca-#1.pdf}}

% ---------- capçalera de la primera pàgina ----------
% #1 matèria (petit), #2 curs, #3 tipus de document i unitat, #4 títol
\newcommand{\capcalera}[4]{%
  \gdef\fm@cursnom{#2}%
  \par\noindent
  \begin{tcolorbox}[colback=curs,colframe=curs,boxrule=0pt,arc=2.2mm,width=3.5cm,height=2.15cm,
      valign=center,left=3mm,right=2mm,top=1mm,bottom=1mm,nobeforeafter,box align=center]
    \color{curstext}{\footnotesize #1}\par\vspace{1pt}{\fontsize{19}{21}\selectfont\bfseries #2\par}%
  \end{tcolorbox}\hspace{5mm}%
  \begin{minipage}[c]{\dimexpr\linewidth-3.5cm-5mm-2.9cm\relax}
    {\color{cursfosc}\small\bfseries\MakeUppercase{#3}}\par\vspace{2pt}
    {\color{navy}\fontsize{20}{23}\selectfont\bfseries #4\par}
  \end{minipage}\hfill
  \begin{minipage}[c]{2.6cm}\raggedleft\includegraphics[width=2.6cm]{\fm@logo}\end{minipage}\par
  \vspace{2pt}}

% ---------- barres de secció del color del curs ----------
\renewcommand{\seccio}[2]{\par\needspace{\fmneed\baselineskip}\bigskip
  \begin{tcolorbox}[colback=curs,colframe=curs,coltext=curstext,arc=2mm,boxrule=0pt,left=3mm,right=3mm,top=1.2mm,bottom=1.2mm]
  {\large\bfseries #1\quad #2}
  \end{tcolorbox}\vspace{-2pt}}

% ---------- estrelles de dificultat del color del curs (a Batxillerat no se'n posen) ----------
\@ifundefined{estrelles}{}{%
  \renewcommand{\estrelles}[1]{{\color{curs}\ifcase#1\or\faStar\faStar[regular]\faStar[regular]\or\faStar\faStar\faStar[regular]\or\faStar\faStar\faStar\fi}}}

% ---------- capçalera i peu de les pàgines ----------
% primera pàgina: logos institucionals al peu
\fancypagestyle{primera}{\fancyhf{}\renewcommand{\headrulewidth}{0pt}
  \fancyfoot[L]{\raisebox{-0.32\height}{\includegraphics[height=0.78cm]{logo-gva.png}}\hspace{5mm}%
    \raisebox{-0.32\height}{\includegraphics[height=0.78cm]{logo-vermellar.png}}}
  \fancyfoot[R]{\scriptsize\color{muted}\href{https://creativecommons.org/licenses/by-nc-sa/4.0/deed.ca}{\ccbyncsa}\ \ CC BY-NC-SA 4.0 · fentmates.org\qquad\small\thepage}}
% la resta: el símbol a l'esquerra i el curs a la dreta, amb el títol al centre com abans
\fancyhead[L]{\raisebox{-0.3\height}{\includegraphics[height=0.62cm]{\fm@marca}}}
\fancyhead[R]{\small\bfseries\color{cursfosc}\fm@cursnom}
\renewcommand{\headrule}{\color{curs!45}\hrule height \headrulewidth}
\makeatother

%   \fmcurs{1eso}
% i, al començament del document:
%   \capcalera{Matemàtiques}{1r ESO}{Fitxa d'exercicis · Unitat 2}{Divisibilitat}
% Opcional abans de \input: \def\fmneed{14} (espai mínim abans de cada secció).
% =====================================================================
\makeatletter
\def\fm@logodir{../../logo/pdf}
\providecommand{\fmneed}{10}
\def\fm@cursnom{}

% clau del curs, color i color del text damunt del color
\newcommand{\fm@curs}[3]{\@namedef{fm@color@#1}{#2}\@namedef{fm@text@#1}{#3}}
\fm@curs{1eso}{E6B84A}{navy}
\fm@curs{2eso}{3B82F6}{white}
\fm@curs{3eso}{FF7A50}{white}
\fm@curs{4eso-a}{A78BFA}{white}
\fm@curs{4eso-b}{6D28D9}{white}
\fm@curs{1batx-mi}{9B1C1C}{white}
\fm@curs{1batx-ccss}{EF6B6B}{white}
\fm@curs{2batx-mii}{1E6B47}{white}
\fm@curs{2batx-ccss}{2BB3A3}{white}

\newcommand{\fmcurs}[1]{%
  \@ifundefined{fm@color@#1}{\PackageError{fentmates}{Curs desconegut: #1}{}}{}%
  \edef\fm@tmp{\noexpand\definecolor{curs}{HTML}{\@nameuse{fm@color@#1}}%
               \noexpand\colorlet{curstext}{\@nameuse{fm@text@#1}}}\fm@tmp
  \colorlet{cursfosc}{curs!72!black}%
  \def\fm@logo{\fm@logodir/vertical-#1.pdf}%
  \def\fm@marca{\fm@logodir/marca-#1.pdf}}

% ---------- capçalera de la primera pàgina ----------
% #1 matèria (petit), #2 curs, #3 tipus de document i unitat, #4 títol
\newcommand{\capcalera}[4]{%
  \gdef\fm@cursnom{#2}%
  \par\noindent
  \begin{tcolorbox}[colback=curs,colframe=curs,boxrule=0pt,arc=2.2mm,width=3.5cm,height=2.15cm,
      valign=center,left=3mm,right=2mm,top=1mm,bottom=1mm,nobeforeafter,box align=center]
    \color{curstext}{\footnotesize #1}\par\vspace{1pt}{\fontsize{19}{21}\selectfont\bfseries #2\par}%
  \end{tcolorbox}\hspace{5mm}%
  \begin{minipage}[c]{\dimexpr\linewidth-3.5cm-5mm-2.9cm\relax}
    {\color{cursfosc}\small\bfseries\MakeUppercase{#3}}\par\vspace{2pt}
    {\color{navy}\fontsize{20}{23}\selectfont\bfseries #4\par}
  \end{minipage}\hfill
  \begin{minipage}[c]{2.6cm}\raggedleft\includegraphics[width=2.6cm]{\fm@logo}\end{minipage}\par
  \vspace{2pt}}

% ---------- barres de secció del color del curs ----------
\renewcommand{\seccio}[2]{\par\needspace{\fmneed\baselineskip}\bigskip
  \begin{tcolorbox}[colback=curs,colframe=curs,coltext=curstext,arc=2mm,boxrule=0pt,left=3mm,right=3mm,top=1.2mm,bottom=1.2mm]
  {\large\bfseries #1\quad #2}
  \end{tcolorbox}\vspace{-2pt}}

% ---------- estrelles de dificultat del color del curs (a Batxillerat no se'n posen) ----------
\@ifundefined{estrelles}{}{%
  \renewcommand{\estrelles}[1]{{\color{curs}\ifcase#1\or\faStar\faStar[regular]\faStar[regular]\or\faStar\faStar\faStar[regular]\or\faStar\faStar\faStar\fi}}}

% ---------- capçalera i peu de les pàgines ----------
% primera pàgina: logos institucionals al peu
\fancypagestyle{primera}{\fancyhf{}\renewcommand{\headrulewidth}{0pt}
  \fancyfoot[L]{\raisebox{-0.32\height}{\includegraphics[height=0.78cm]{logo-gva.png}}\hspace{5mm}%
    \raisebox{-0.32\height}{\includegraphics[height=0.78cm]{logo-vermellar.png}}}
  \fancyfoot[R]{\scriptsize\color{muted}\href{https://creativecommons.org/licenses/by-nc-sa/4.0/deed.ca}{\ccbyncsa}\ \ CC BY-NC-SA 4.0 · fentmates.org\qquad\small\thepage}}
% la resta: el símbol a l'esquerra i el curs a la dreta, amb el títol al centre com abans
\fancyhead[L]{\raisebox{-0.3\height}{\includegraphics[height=0.62cm]{\fm@marca}}}
\fancyhead[R]{\small\bfseries\color{cursfosc}\fm@cursnom}
\renewcommand{\headrule}{\color{curs!45}\hrule height \headrulewidth}
\makeatother

%   \fmcurs{1eso}
% i, al començament del document:
%   \capcalera{Matemàtiques}{1r ESO}{Fitxa d'exercicis · Unitat 2}{Divisibilitat}
% Opcional abans de \input: \def\fmneed{14} (espai mínim abans de cada secció).
% =====================================================================
\makeatletter
\def\fm@logodir{../../logo/pdf}
\providecommand{\fmneed}{10}
\def\fm@cursnom{}

% clau del curs, color i color del text damunt del color
\newcommand{\fm@curs}[3]{\@namedef{fm@color@#1}{#2}\@namedef{fm@text@#1}{#3}}
\fm@curs{1eso}{E6B84A}{navy}
\fm@curs{2eso}{3B82F6}{white}
\fm@curs{3eso}{FF7A50}{white}
\fm@curs{4eso-a}{A78BFA}{white}
\fm@curs{4eso-b}{6D28D9}{white}
\fm@curs{1batx-mi}{9B1C1C}{white}
\fm@curs{1batx-ccss}{EF6B6B}{white}
\fm@curs{2batx-mii}{1E6B47}{white}
\fm@curs{2batx-ccss}{2BB3A3}{white}

\newcommand{\fmcurs}[1]{%
  \@ifundefined{fm@color@#1}{\PackageError{fentmates}{Curs desconegut: #1}{}}{}%
  \edef\fm@tmp{\noexpand\definecolor{curs}{HTML}{\@nameuse{fm@color@#1}}%
               \noexpand\colorlet{curstext}{\@nameuse{fm@text@#1}}}\fm@tmp
  \colorlet{cursfosc}{curs!72!black}%
  \def\fm@logo{\fm@logodir/vertical-#1.pdf}%
  \def\fm@marca{\fm@logodir/marca-#1.pdf}}

% ---------- capçalera de la primera pàgina ----------
% #1 matèria (petit), #2 curs, #3 tipus de document i unitat, #4 títol
\newcommand{\capcalera}[4]{%
  \gdef\fm@cursnom{#2}%
  \par\noindent
  \begin{tcolorbox}[colback=curs,colframe=curs,boxrule=0pt,arc=2.2mm,width=3.5cm,height=2.15cm,
      valign=center,left=3mm,right=2mm,top=1mm,bottom=1mm,nobeforeafter,box align=center]
    \color{curstext}{\footnotesize #1}\par\vspace{1pt}{\fontsize{19}{21}\selectfont\bfseries #2\par}%
  \end{tcolorbox}\hspace{5mm}%
  \begin{minipage}[c]{\dimexpr\linewidth-3.5cm-5mm-2.9cm\relax}
    {\color{cursfosc}\small\bfseries\MakeUppercase{#3}}\par\vspace{2pt}
    {\color{navy}\fontsize{20}{23}\selectfont\bfseries #4\par}
  \end{minipage}\hfill
  \begin{minipage}[c]{2.6cm}\raggedleft\includegraphics[width=2.6cm]{\fm@logo}\end{minipage}\par
  \vspace{2pt}}

% ---------- barres de secció del color del curs ----------
\renewcommand{\seccio}[2]{\par\needspace{\fmneed\baselineskip}\bigskip
  \begin{tcolorbox}[colback=curs,colframe=curs,coltext=curstext,arc=2mm,boxrule=0pt,left=3mm,right=3mm,top=1.2mm,bottom=1.2mm]
  {\large\bfseries #1\quad #2}
  \end{tcolorbox}\vspace{-2pt}}

% ---------- estrelles de dificultat del color del curs (a Batxillerat no se'n posen) ----------
\@ifundefined{estrelles}{}{%
  \renewcommand{\estrelles}[1]{{\color{curs}\ifcase#1\or\faStar\faStar[regular]\faStar[regular]\or\faStar\faStar\faStar[regular]\or\faStar\faStar\faStar\fi}}}

% ---------- capçalera i peu de les pàgines ----------
% primera pàgina: logos institucionals al peu
\fancypagestyle{primera}{\fancyhf{}\renewcommand{\headrulewidth}{0pt}
  \fancyfoot[L]{\raisebox{-0.32\height}{\includegraphics[height=0.78cm]{logo-gva.png}}\hspace{5mm}%
    \raisebox{-0.32\height}{\includegraphics[height=0.78cm]{logo-vermellar.png}}}
  \fancyfoot[R]{\scriptsize\color{muted}\href{https://creativecommons.org/licenses/by-nc-sa/4.0/deed.ca}{\ccbyncsa}\ \ CC BY-NC-SA 4.0 · fentmates.org\qquad\small\thepage}}
% la resta: el símbol a l'esquerra i el curs a la dreta, amb el títol al centre com abans
\fancyhead[L]{\raisebox{-0.3\height}{\includegraphics[height=0.62cm]{\fm@marca}}}
\fancyhead[R]{\small\bfseries\color{cursfosc}\fm@cursnom}
\renewcommand{\headrule}{\color{curs!45}\hrule height \headrulewidth}
\makeatother

\fmcurs{4eso-a}

\begin{document}
\thispagestyle{primera}

\capcalera{Matemáticas A}{4\textsuperscript{o} ESO}{Ficha de ejercicios · Unidad 3}{\raggedright Polinomios}
\vspace{10pt}
Nombre y apellidos: \rule{6.6cm}{0.4pt}\hfill Grupo: \rule{1.6cm}{0.4pt}\hfill Fecha: \rule{2.2cm}{0.4pt}
\vspace{8pt}

\begin{tcolorbox}[colback=soft,colframe=soft,arc=2mm,left=3mm,right=3mm,top=1.5mm,bottom=1.5mm]
\small Ejercicios graduados por niveles de dificultad:\quad \estrelles{1}\ básico \qquad \estrelles{2}\ medio \qquad \estrelles{3}\ avanzado.\\[2pt] Resuélvelos en el cuaderno mostrando todos los pasos. Puedes comprobar los resultados con las herramientas interactivas de los apuntes.
\end{tcolorbox}

% =====================================================
\seccio{1}{Monomios}

\exf{1} Completa la tabla:\begin{center}\renewcommand{\arraystretch}{1.5}\begin{tabular}{|C{2.4cm}|C{2.6cm}|C{2.6cm}|C{1.6cm}|}\hline\rowcolor{soft}\textbf{Monomio} & \textbf{Coeficiente} & \textbf{Parte literal} & \textbf{Grado}\\\hline$-3x^4$ &  &  & \\\hline$x^2y$ &  &  & \\\hline$\frac25x$ &  &  & \\\hline$-a^3b^2$ &  &  & \\\hline$8$ &  &  & \\\hline$-x$ &  &  & \\\hline\end{tabular}\end{center}

\ex{1} Agrupa los monomios semejantes: \ $3x^2$, \ $-5x$, \ $x^3$, \ $7x^2$, \ $2x$, \ $-x^3$, \ $\frac12x^2$.

\ex{1} Reduce:\begin{tasks}(3)
  \task $5x+3x-x$
  \task $4x^2-7x^2+x^2$
  \task $2x^3+5x-x^3+3x$
  \task $3x^2-2x+5-x^2+4x-8$
  \task $\frac12x+\frac34x$
  \task $6xy-2xy+xy$
\end{tasks}\recorda{Solo se suman o restan los monomios semejantes: se suman los coeficientes y se deja la misma parte literal.}

\ex{1} Calcula:\begin{tasks}(3)
  \task $3x^2\cdot4x^3$
  \task $(-2x)\cdot5x^4$
  \task $(-x^3)\cdot(-6x^2)$
  \task $\frac23x\cdot\frac94x^2$
  \task $4x^2y\cdot3xy^3$
  \task $2x\cdot(-3x)\cdot x^2$
\end{tasks}\recorda{En el producto, se multiplican los coeficientes y se suman los exponentes de cada letra: $x^2\cdot x^3=x^5$.}

\ex{2} Calcula:\begin{tasks}(3)
  \task $12x^5:3x^2$
  \task $-20x^6:4x$
  \task $9x^4:(-3x^4)$
  \task $(2x^3)^2$
  \task $(-3x^2)^3$
  \task $\left(\frac12x\right)^4$
\end{tasks}

\ex{2} Calcula el valor numérico de cada monomio para los valores indicados:\begin{tasks}(1)
  \task $-2x^3$ \ para \ $x=-1,\ 2,\ \frac{1}{2}$
  \task $\frac34x^2$ \ para \ $x=-2,\ 0,\ 4$
\end{tasks}


% =====================================================
\seccio{2}{Polinomios y valor numérico}

\exf{1} Completa la tabla:\begin{center}\renewcommand{\arraystretch}{1.5}\small\begin{tabular}{|l|C{3.7cm}|C{1.1cm}|C{1.5cm}|C{1.6cm}|C{1.95cm}|}\hline\rowcolor{soft}\textbf{Polinomio} & \textbf{Ordenado} & \textbf{Grado} & \textbf{Térm. indep.} & \textbf{Coef. princ.} & \textbf{¿Completo?}\\\hline$P(x)=3x^2-5x^4+x-7$ &  &  &  &  & \\\hline$Q(x)=6-2x+x^3$ &  &  &  &  & \\\hline$R(x)=4x-x^2+1$ &  &  &  &  & \\\hline$S(x)=2x^5+x^2-3x$ &  &  &  &  & \\\hline$U(x)=\frac12x^3-x^2+3x-2$ &  &  &  &  & \\\hline\end{tabular}\end{center}

\ex{1} Escribe, en cada caso, un polinomio que cumpla las condiciones:\begin{tasks}(2)
  \task de grado 3, completo y con término independiente $-2$;
  \task de grado 4 y solo dos términos;
  \task de grado 2, con coeficiente principal $-3$;
  \task de grado 5 y sin término independiente.
\end{tasks}

\ex{1} Dado el polinomio $P(x)=x^3-2x^2+4x-1$, calcula $P(0)$, $P(1)$, $P(-1)$, $P(2)$ y $P(-2)$.\recorda{Sustituimos la $x$ por el número, entre paréntesis si es negativo: $(-2)^3=-8$, $(-2)^2=4$.}

\ex{1} Calcula el valor numérico de los polinomios siguientes para $x=3$, $x=-1$ y $x=0$:\begin{tasks}(3)
  \task $P(x)=2x^{2}-5x+3$
  \task $P(x)=-x^{3}+4x$
  \task $P(x)=x^{4}-3x^{2}+2$
\end{tasks}

\ex{2} ¿Qué valor tiene que tener $k$ para que el polinomio $P(x)=2x^3-kx+5$ valga 9 para $x=2$?

\ex{2} Halla $a$ en el polinomio $P(x)=x^2+ax-4$ si sabemos que $P(-1)=-6$.

\ex{2} El coste, en euros, de fabricar $x$ camisetas es $C(x)=0,02x^2+3x+50$. ¿Cuánto cuesta fabricar 10? ¿Y 100?

\ex{3} La altura (en metros) de una pelota lanzada hacia arriba es $h(t)=-5t^2+20t+1$, donde $t$ es el tiempo en segundos. Calcula $h(0)$, $h(1)$, $h(2)$ y $h(4)$. ¿Qué significan $h(0)$ y $h(4)$?


% =====================================================
\seccio{3}{Suma y resta de polinomios}

\exf{1} Dados $P(x)=3x^3-2x^2+x-5$, $Q(x)=x^3+4x^2-6$ y $R(x)=-2x^2+3x+1$, calcula:\begin{tasks}(2)
  \task $P(x)+Q(x)$
  \task $P(x)-Q(x)$
  \task $Q(x)-R(x)$
  \task $P(x)+Q(x)+R(x)$
\end{tasks}\recorda{Para restar, cambiamos el signo de \emph{todos} los términos del polinomio que restamos: $-(x^2-3x)=-x^2+3x$.}

\ex{2} Con los mismos polinomios, calcula:\begin{tasks}(3)
  \task $P(x)-[Q(x)+R(x)]$
  \task $2P(x)-3R(x)$
  \task $P(x)-2Q(x)+R(x)$
\end{tasks}

\ex{2} Halla el polinomio $Q(x)$ tal que $\left(x^3-2x+4\right)+Q(x)=3x^3+x^2-5$.

\ex{2} Los lados de un triángulo miden $x+3$, $2x-1$ y $x^2+x$ centímetros. Escribe el polinomio que da el perímetro y calcúlalo para $x=4$.


% =====================================================
\seccio{4}{Producto de polinomios}

\exf{1} Multiplica:\begin{tasks}(2)
  \task $3x\left(x^2-2x+5\right)$
  \task $-2x^2\left(x^3-4x+1\right)$
  \task $\frac12x\left(4x^2-6x+2\right)$
  \task $5x^3\left(-x^2+x-3\right)$
  \task $\left(x^2-3x+1\right)\cdot(-4x)$
  \task $-x\left(2x^4-x^2+7\right)$
\end{tasks}\recorda{Se multiplica el monomio por cada término del polinomio. Cuidado con los signos: $(-)\cdot(-)=+$.}

\ex{1} Multiplica y reduce:\begin{tasks}(2)
  \task $(x+2)(x-5)$
  \task $(2x-1)(3x+4)$
  \task $\left(x^2+1\right)(x-3)$
  \task $\left(x^2-2x+3\right)(x+1)$
  \task $\left(2x^2-x+4\right)(x-2)$
  \task $\left(x^2+x-1\right)\left(x^2-x+2\right)$
\end{tasks}\recorda{Cada término del primer polinomio se multiplica por cada término del segundo; después se reducen los términos semejantes.}

\ex{2} Dados $P(x)=x^2-3x+2$, $Q(x)=2x+1$ y $R(x)=x-4$, calcula:\begin{tasks}(3)
  \task $P(x)\cdot Q(x)$
  \task $P(x)\cdot R(x)-Q(x)$
  \task $[Q(x)+R(x)]\cdot P(x)$
\end{tasks}

\ex{2} Opera y reduce:\begin{tasks}(1)
  \task $2x(x-3)-x(x+1)$
  \task $(x+2)(x-1)-(x-3)(x+4)$
  \task $3\left(x^2-1\right)-(x+1)(2x-3)$
\end{tasks}

\ex{3} Un rectángulo mide $2x+3$ cm de base y $x-1$ cm de altura. Escribe el polinomio que da el área y calcúlala para $x=5$. ¿Y el volumen de una caja de aristas $x$, $x+2$ y $x+5$? Calcúlalo para $x=3$.


% =====================================================
\seccio{5}{Identidades notables}

\exf{1} Desarrolla aplicando las identidades notables:\begin{tasks}(3)
  \task $(x+3)^2$
  \task $(x-4)^2$
  \task $(x+5)(x-5)$
  \task $(2x+1)^2$
  \task $(3x-2)^2$
  \task $(2x+7)(2x-7)$
  \task $\left(x^2+3\right)^2$
  \task $(5-x)^2$
  \task $\left(\frac12x+1\right)^2$
\end{tasks}\recorda{$(a+b)^2=a^2+2ab+b^2$, \quad $(a-b)^2=a^2-2ab+b^2$, \quad $(a+b)(a-b)=a^2-b^2$.}

\ex{2} Opera y simplifica:\begin{tasks}(2)
  \task $(x+1)^2-(x-1)^2$
  \task $(x+3)^2-(x+3)(x-3)$
  \task $(2x-1)^2+(x+2)(x-2)$
  \task $2(x-3)^2-(x+1)^2$
\end{tasks}

\ex{2} Completa los huecos:\begin{tasks}(1)
  \task $(x+\square)^2=x^2+\square+36$
  \task $(\square-5)^2=4x^2-\square+25$
  \task $(3x+\square)(3x-\square)=\square-16$
\end{tasks}

\ex{3} Calcula sin calculadora, usando las identidades notables: \ $101^2$, \ $99^2$ \ y \ $51\cdot49$.


% =====================================================
\seccio{6}{División de polinomios}

\exf{1} Divide:\begin{tasks}(1)
  \task $\left(6x^3-9x^2+3x\right):3x$
  \task $\left(10x^5-15x^3+5x^2\right):5x^2$
  \task $\left(8x^4-4x^3+2x^2\right):\left(-2x^2\right)$
\end{tasks}

\ex{2} Haz las divisiones siguientes e indica el cociente y el resto:\begin{tasks}(1)
  \task $\left(2x^3+3x^2-x+4\right):\left(x^2+1\right)$
  \task $\left(x^4-2x^3+3x-1\right):\left(x^2-x\right)$
  \task $\left(6x^3-x^2+x+5\right):(2x-1)$
  \task $\left(x^4-1\right):\left(x^2+1\right)$
  \task $\left(3x^4+2x^2-x\right):\left(x^2-2x+1\right)$
\end{tasks}\recorda{Ordena el dividendo y deja hueco para los términos que faltan. Para cuando el resto tenga grado menor que el divisor.}

\ex{2} Comprueba el resultado de la división a) del ejercicio anterior con la prueba de la división, $D(x)=d(x)\cdot c(x)+r(x)$.

\ex{3} En una división, el divisor es $x^2-2$, el cociente es $x+3$ y el resto es $2x-1$. ¿Cuál es el dividendo?

\ex{3} El área de un rectángulo es $2x^2+7x+3$ y la base mide $x+3$. ¿Cuánto mide la altura? ¿Cuál es el perímetro?


% =====================================================
\seccio{7}{Regla de Ruffini}

\exf{1} Divide aplicando la regla de Ruffini e indica el cociente y el resto:\begin{tasks}(2)
  \task $\left(x^{3}-4x^{2}+x+6\right):(x-2)$
  \task $\left(x^{3}+2x^{2}-x+3\right):(x+1)$
  \task $\left(2x^{3}-3x^{2}+4\right):(x-1)$
  \task $\left(x^{4}-3x^{2}+2x-5\right):(x+2)$
  \task $\left(x^{4}-16\right):(x-2)$
  \task $\left(3x^{3}-5x+1\right):(x+3)$
\end{tasks}\recorda{Escribe todos los coeficientes, con un 0 por cada término que falte. Si el divisor es $x+2$, a la izquierda va el $-2$.}

\ex{2} Divide por Ruffini:\begin{tasks}(2)
  \task $\left(x^5-1\right):(x-1)$
  \task $\left(x^3+27\right):(x+3)$
\end{tasks}

\ex{2} Usa la regla de Ruffini para saber si $x^3-2x^2-5x+6$ es divisible entre $x-3$ y entre $x+1$.

\ex{3} Halla el valor de $k$ para que la división $\left(x^3+kx^2-4x+12\right):(x-2)$ sea exacta.


% =====================================================
\seccio{8}{Valor numérico con Ruffini. Teorema del resto}

\exf{1} Calcula, sin hacer la división, el resto de:\begin{tasks}(2)
  \task $\left(x^{3}-5x+2\right):(x-2)$
  \task $\left(x^{4}+2x^{3}-x\right):(x+1)$
  \task $\left(2x^{5}-3x^{2}+1\right):(x-1)$
  \task $\left(x^{3}+x^{2}+x+1\right):(x+2)$
\end{tasks}\recorda{El resto de la división $P(x):(x-a)$ es el valor numérico $P(a)$.}

\ex{2} Dado $P(x)=x^4-2x^3-5x^2+x+3$, calcula $P(3)$ y $P(-2)$ con la regla de Ruffini. Compruébalo sustituyendo.

\ex{2} Halla $m$ para que el resto de la división $\left(x^3-mx+4\right):(x+2)$ sea 10.

\ex{3} ¿Qué valor tiene que tener $a$ para que el polinomio $2x^3+ax^2-x+3$ dé resto 5 al dividirlo entre $x-1$?


% =====================================================
\seccio{9}{Raíces de un polinomio}

\exf{1} Comprueba si $x=1$, $x=-1$ y $x=2$ son raíces del polinomio $P(x)=x^3-2x^2-x+2$.

\ex{1} ¿Cuáles de estos números son raíces de $x^3+x^2-6x$? \ $0$, \ $2$, \ $-2$, \ $-3$.

\ex{2} Halla las raíces enteras de los polinomios:\begin{tasks}(2)
  \task $x^{3}-6x^{2}+11x-6$
  \task $x^{3}+2x^{2}-x-2$
  \task $x^{3}-3x+2$
  \task $x^{4}-5x^{2}+4$
\end{tasks}\recorda{Las raíces enteras son divisores del término independiente: pruébalos con Ruffini (o calculando el valor numérico).}

\ex{2} Halla las raíces de los polinomios de segundo grado resolviendo la ecuación:\begin{tasks}(2)
  \task $x^{2}-7x+12$
  \task $2x^{2}-x-1$
  \task $x^{2}+2x+5$
  \task $x^{2}-6x+9$
\end{tasks}

\ex{2} Sin hacer ningún cálculo, explica por qué $x=5$ no puede ser una raíz de $x^3+x^2-3x+6$.

\ex{3} Escribe un polinomio de grado 3, con coeficiente principal 1, que tenga por raíces $2$, $-1$ y $3$. Halla también $k$ para que $x=2$ sea raíz de $x^3+kx-6$.


% =====================================================
\seccio{10}{Factor común e identidades notables}

\exf{1} Saca factor común:\begin{tasks}(3)
  \task $4x^{2}-8x$
  \task $3x^{3}+6x^{2}-9x$
  \task $5x^{4}-10x^{3}$
  \task $12x^{5}-18x^{3}+6x^{2}$
  \task $x^{3}+x^{2}+x$
  \task $2x^{4}-14x^{2}$
\end{tasks}\recorda{Se saca el m.c.d. de los coeficientes y la menor potencia de $x$: $6x^3-9x^2=3x^2(2x-3)$.}

\ex{1} Factoriza usando las identidades notables:\begin{tasks}(3)
  \task $x^{2}-36$
  \task $x^{2}+8x+16$
  \task $x^{2}-12x+36$
  \task $9x^{2}-4$
  \task $4x^{2}+20x+25$
  \task $25x^{2}-10x+1$
  \task $x^{4}-9$
  \task $x^{2}-\frac{1}{4}$
  \task $16-x^2$
\end{tasks}\recorda{$a^2+2ab+b^2=(a+b)^2$, \quad $a^2-2ab+b^2=(a-b)^2$, \quad $a^2-b^2=(a+b)(a-b)$.}

\ex{2} Factoriza: primero saca factor común y después aplica una identidad notable.\begin{tasks}(3)
  \task $2x^{3}-18x$
  \task $3x^{3}+12x^{2}+12x$
  \task $x^{4}-x^{2}$
  \task $5x^{3}-10x^{2}+5x$
  \task $x^{5}-16x$
  \task $2x^{4}-50x^{2}$
\end{tasks}

\ex{3} Calcula sin calculadora, factorizando: \ $83^2-17^2$.


% =====================================================
\seccio{11}{Factorización de polinomios}

\exf{2} Factoriza estos polinomios de segundo grado (busca primero sus raíces):\begin{tasks}(2)
  \task $x^{2}-5x+6$
  \task $x^{2}+x-12$
  \task $2x^{2}-8x+6$
  \task $x^{2}-4x-5$
\end{tasks}\recorda{Si $x_1$ y $x_2$ son las raíces de $ax^2+bx+c$, entonces $ax^2+bx+c=a(x-x_1)(x-x_2)$.}

\ex{2} Factoriza con la regla de Ruffini:\begin{tasks}(2)
  \task $x^{3}-2x^{2}-5x+6$
  \task $x^{3}+4x^{2}+x-6$
  \task $x^{3}-7x-6$
  \task $x^{3}-3x^{2}-4x+12$
\end{tasks}

\ex{3} Factoriza (si se puede, saca primero factor común):\begin{tasks}(2)
  \task $x^{4}-x^{3}-7x^{2}+x+6$
  \task $2x^{3}-4x^{2}-2x+4$
  \task $x^{4}-2x^{3}-3x^{2}$
  \task $x^{3}-x^{2}-8x+12$
\end{tasks}

\ex{3} Factoriza e indica las raíces y cuántas veces se repite cada una:\begin{tasks}(2)
  \task $x^{3}-3x^{2}+3x-1$
  \task $x^{4}-8x^{2}+16$
\end{tasks}


% =====================================================
\seccio{12}{Problemas}

\exf{2} A un cuadrado de lado $x+2$ le recortamos un cuadrado de lado $x$. Escribe el polinomio que da el área que queda y calcúlala para $x=3$ cm.

\ex{2} De una lámina de cartón de $20\times30$ cm recortamos un cuadrado de lado $x$ en cada esquina y doblamos los bordes para hacer una caja sin tapa. Escribe el polinomio $V(x)$ del volumen de la caja y calcúlalo para $x=5$ cm.

\ex{3} El área de un rectángulo es $x^2+5x+6$ y uno de los lados mide $x+2$. ¿Cuánto mide el otro lado? Escribe el perímetro y calcúlalo para $x=4$.

\ex{3} Piensa un número, súmale 3, eleva el resultado al cuadrado y réstale el cuadrado del número pensado. Comprueba con polinomios que el resultado siempre es el séxtuplo del número más 9.

\vfill
\begin{center}\color{muted}\small\faCheckCircle[regular]\ \ ¡Ánimo! Comprueba tus resultados con las herramientas interactivas de los apuntes de la unidad.\end{center}

\end{document}
